\documentclass[a4paper]{article}

\usepackage[margin=1in]{geometry}
\usepackage{amsmath,amsfonts,mathtools,amsthm,amssymb}
\usepackage[l2tabu,orthodox]{nag}
\usepackage{microtype} % fixes spacing or whatever
\usepackage{enumitem}
\usepackage{tikz-cd}
\usepackage{xcolor}
\usepackage{physics}
\usepackage{mathrsfs}
\usepackage{cancel}
% \usepackage{libertine,libertinust1math}
\usepackage{eucal}
\usepackage[toc,page]{appendix}
\usepackage{pgfplots}
\pgfplotsset{compat=1.18}

\newtheorem{proposition}{Proposition}
\newtheorem{theorem}{Theorem}
\newtheorem{lemma}{Lemma}
\newtheorem{corollary}{Corollary}

\theoremstyle{definition}
\newtheorem{remark}{Remark}
\newtheorem{definition}{Definition}
\newtheorem{example}{Example}

% \usepackage[sc,osf]{mathpazo}   % With old-style figures and real smallcaps.
% \linespread{1.025}              % Palatino leads a little more leading
% % Euler for math and numbers
% \usepackage[euler-digits,small]{eulervm}
% \AtBeginDocument{\renewcommand{\hbar}{\hslash}}

\usepackage{etoolbox}
\usepackage{dynkin-diagrams}
\def\row#1/#2!{\dynkin{#1}{#2}\\}
\newcommand{\tble}[1]{
   \renewcommand*\do[1]{\row##1!}
   \[
      \begin{array}{ll}\docsvlist{#1}\end{array}
   \]
}

\allowdisplaybreaks

\renewcommand{\epsilon}{\varepsilon}
% \renewcommand{\phi}{\varphi}

\DeclareMathAlphabet{\mathpzc}{OT1}{pzc}{m}{it}

\newcommand{\goth}[1]{\mathfrak{#1}}
\newcommand{\red}[1]{#1_{\text{red}}}
\newcommand{\ad}[1]{#1_{\text{Ad}}}
\newcommand{\reg}[1]{#1_{\text{reg}}}
\newcommand{\sh}[1]{#1^{\text{sh}}}
\newcommand{\cpdv}[2]{\cfrac{\partial #1}{\partial #2}}

\newcommand{\fA}{\mathcal{A}}
\newcommand{\fB}{\mathcal{B}}
\newcommand{\fC}{\mathcal{C}}
\newcommand{\fD}{\mathcal{D}}
\newcommand{\fE}{\mathcal{E}}
\newcommand{\fF}{\mathcal{F}}
\newcommand{\fG}{\mathcal{G}}
\newcommand{\fH}{\mathcal{H}}
\newcommand{\fI}{\mathcal{I}}
\newcommand{\fJ}{\mathcal{J}}
\newcommand{\fK}{\mathcal{K}}
\newcommand{\fL}{\mathcal{L}}
\newcommand{\fM}{\mathcal{M}}
\newcommand{\fN}{\mathcal{N}}
\newcommand{\fO}{\mathcal{O}}
\newcommand{\fP}{\mathcal{P}}
\newcommand{\fQ}{\mathcal{Q}}
\newcommand{\fR}{\mathcal{R}}
\newcommand{\fS}{\mathcal{S}}
\newcommand{\fT}{\mathcal{T}}
\newcommand{\fX}{\mathcal{X}}
\newcommand{\fY}{\mathcal{Y}}
\newcommand{\fZ}{\mathcal{Z}}

\newcommand{\be}{\mathbf{e}}
\newcommand{\bx}{\mathbf{x}}
\newcommand{\bone}{\mathbf{1}}
\newcommand{\cx}{\bar{x}}
\newcommand{\cy}{\bar{y}}
\newcommand{\vx}{\vec{x}}
\newcommand{\vsigma}{\vec{\sigma}}
\newcommand{\tzeta}{\tilde{\zeta}}

\newcommand{\gm}{\goth{m}}
\newcommand{\gp}{\goth{p}}
\newcommand{\gF}{\goth{F}}
\newcommand{\gU}{\goth{U}}
\newcommand{\gV}{\goth{V}}

\newcommand{\A}{\mathbf{A}}
\newcommand{\C}{\mathbf{C}}
\newcommand{\F}{\mathbf{F}}
\newcommand{\G}{\mathbf{G}}
\newcommand{\bH}{\mathbf{H}}
\newcommand{\N}{\mathbf{N}}
\newcommand{\PP}{\mathbf{P}}
\newcommand{\Q}{\mathbf{Q}}
\newcommand{\R}{\mathbf{R}}
\newcommand{\Z}{\mathbf{Z}}

\newcommand{\dlog}{\dd\log}

\newcommand\srestr[2]{{\left.\kern-\nulldelimiterspace #1\vphantom{\small|} \right|_{#2}}}
\newcommand\restr[2]{{\left.\kern-\nulldelimiterspace #1 \vphantom{\big|} \right|_{#2}}}
\newcommand\gsec[2]{\Gamma({#1}, {#2})}
\newcommand\gsecx[1]{\Gamma(X, {#1})}

\DeclareMathOperator{\chr}{char}
\DeclareMathOperator{\rProj}{\mathbf{Proj}}
\DeclareMathOperator{\rSpec}{\mathbf{Spec}}
\DeclareMathOperator{\rHom}{\mathbf{Hom}}
\DeclareMathOperator{\rExt}{\mathbf{Ext}}
\DeclareMathOperator{\id}{id}
\DeclareMathOperator{\Frac}{Frac}
\DeclareMathOperator{\rk}{rank}
\DeclareMathOperator{\deter}{det}
\DeclareMathOperator{\pic}{Pic}
\DeclareMathOperator{\cacl}{CaCl}
\DeclareMathOperator{\trd}{tr.d.}
\DeclareMathOperator{\cl}{Cl}
\DeclareMathOperator{\depth}{depth}
\DeclareMathOperator{\codim}{codim}
\DeclareMathOperator{\Div}{Div}
\DeclareMathOperator{\coker}{coker}
\DeclareMathOperator{\len}{length}
\DeclareMathOperator{\height}{height}
\DeclareMathOperator{\supp}{Supp}
\DeclareMathOperator{\proj}{Proj}
\DeclareMathOperator{\im}{im}
\DeclareMathOperator{\Hom}{Hom}
\DeclareMathOperator{\Der}{Der}
\DeclareMathOperator{\spec}{Spec}
\DeclareMathOperator{\Aut}{Aut}
\DeclareMathOperator{\ch}{char}
\DeclareMathOperator{\tor}{Tor}
\DeclareMathOperator{\Ann}{Ann}
\DeclareMathOperator{\Syl}{Syl}
\DeclareMathOperator{\Sym}{Sym}
\DeclareMathOperator{\GL}{GL}
\DeclareMathOperator{\SL}{SL}
\DeclareMathOperator{\Stab}{Stab}
\DeclareMathOperator{\Perm}{Perm}
\DeclareMathOperator{\Orb}{Orb}
\DeclareMathOperator{\Gal}{Gal}
\DeclareMathOperator{\Supp}{Supp}
\DeclareMathOperator{\Ext}{Ext}
\DeclareMathOperator{\Tor}{Tor}
\DeclareMathOperator{\PGL}{PGL}
\DeclareMathOperator{\hd}{hd}
\DeclareMathOperator{\pd}{pd}
\DeclareMathOperator{\SO}{SO}
\DeclareMathOperator{\SU}{SU}
\DeclareMathOperator{\Sp}{Sp}
\DeclareMathOperator{\Oth}{O}
\DeclareMathOperator{\Uni}{U}
\DeclareMathOperator{\evf}{ev}
\DeclareMathOperator{\cH}{H}
\DeclareMathOperator{\dR}{R}
\DeclareMathOperator{\End}{End}
\DeclareMathOperator{\Hasse}{Hasse}
\DeclareMathOperator{\Num}{Num}

\author{James Lee}

% Solarized
% \pagecolor[RGB]{0,20,26}
% \color[RGB]{191,191,191}

% Sephia
% \pagecolor[RGB]{249,239,220}

% Grey
% \pagecolor[RGB]{200, 200, 200}

% Black
% \pagecolor[RGB]{0,0,0}
% \color[RGB]{255,255,255}


\title{Chapter 4, Section 4}

\begin{document}

\maketitle

\begin{enumerate} [label=\textbf{\arabic*.}, leftmargin=0em]

    \item[\textbf{3.}] Let the elliptic curve $X$ be embedded in $\P^2$ so as to have the equation $y^2 = x(x - 1)(x - \lambda)$.
          Show that any automorphism of $X$ leaving $P_0 = (0:1:0)$ fixed is induced by an automorphism of $\P^2$ coming from the automorphism of the affine $(x, y)$-plane given by
          \begin{equation*}
              \begin{cases*}
                  x' = ax + b \\
                  y' = cy.
              \end{cases*}
          \end{equation*}
          In each of the four cases of (4.7), describe these automorphisms of $\P^2$ explicitly, and hence determine the structure of the group $G = \Aut(X, P_0)$.

          \begin{proof} $ $ \vspace{0pt}
              Let $ : X \to \P^1$ be the unique $g_2^1$ with $f(P_0) = \infty$, branched over $0, 1, \lambda, \infty$.
              If $\sigma \in G$, then by (4.4) there exists an automorphism $\tau$ of $\P^1$ which fixes $\infty$ and satisfies $f \circ \sigma = \tau \circ f$.
              In particular, $\tau$ sends $\{0, 1, \lambda\}$ to $\{0, 1, \lambda\}$ in some order.
              If $\tau = \id$, then either $\sigma = \id$ or $\sigma$ is the unique automorphism of order 2 fixing $P_0$ and interchanging the sheets of $f$.
              We call it the \textit{hyperbolic involution} of $X$, and we denote it by $\iota : X \to X$.
              It comes from the automorphism of the affine $(x, y)$-plane corresponding to the reflection across the $x$-axis:
              \begin{equation*}
                  \begin{cases*}
                      x' = x \\
                      y' = -y
                  \end{cases*}.
              \end{equation*}
              The other automorphims are determined using the methods of (4.7).
              In particular, we look at the orbits of the action of $S_3$ on $k - \{0, 1\}$ (4.5).

              \fbox{$j \neq 0, 1728$} Since the orbits of $\lambda$ are all distinct, $\tau = \id$, and $\sigma = \id$ or $\iota$.

              \fbox{$j = 1728, \car{k} \neq 3$}
              This is the case $\lambda = -1, \frac{1}{2},$ or $2$, and $\lambda$ coincides with one other element of its orbit under the action of $S_3$.
              For example, if $\lambda = -1$, then $\lambda = 1/\lambda$, and $G$ is the cyclic group of order 4 generated by
              \begin{equation*}
                  \begin{cases*}
                      x' = - x \\
                      y' = iy
                  \end{cases*}.
              \end{equation*}

              \fbox{$j = 0, \car{k} \neq 3$}
              This is the case $\lambda = -\omega$ or $-\omega^2$, where $\omega$ is a third root of unity, and $\lambda$ coincides with two other elements of its orbit under $S_3$.
              For example, if $\lambda = -\omega$, then
              \begin{equation*}
                  \lambda = \frac{1}{1 - \lambda} = \frac{\lambda - 1}{\lambda},
              \end{equation*}
              and $G$ is the cyclic group of order 6 generated by
              \begin{equation*}
                  \begin{cases*}
                      x' = \omega^2 x + 1 \\
                      y' = -y
                  \end{cases*}.
              \end{equation*}

              \fbox{$j = 0 ~ (= 1728), \car{k} = 3$}
              If $\car{k} = 3$ and $j = 0$, then $\lambda = -1$, and all six elements of its orbit under $S_3$ coincide.
              Note that $k$ has no nontrivial third roots of unity, and the values $\lambda = -1, \frac{1}{2}, 2$ of the previous case coincide too.
              So $G$ is the cyclic group of order 12 generated by
              \begin{equation*}
                  \begin{cases*}
                      x' = x + 1 \\
                      y' = iy
                  \end{cases*}.
              \end{equation*}
          \end{proof}

    \item[\textbf{4.}] Let $X$ be an elliptic curve in $\P^2$ given by an equation of the form
          \begin{equation*}
              y^2 + a_1xy + a_3y = x^3 + a_2x^2 + a_4x + a_6.
          \end{equation*}
          Show that the $j$-invariant is a rational function of the $a_i$, with coefficients in $\Q$.
          In particular, if the $a_i$ are all in some field $k_0 \subseteq k$, then $j \in k_0$ also.
          Furthermore, for every $\alpha \in k_0$, there exists an elliptic curve defined over $k_0$ with $j$-invariant equal to $\alpha$.

          \begin{proof}
              Complete the square in $y$ to get an equation of the form
              \begin{equation*}
                  y^2 = x^3 + Ax^2 + Bx + C,
              \end{equation*}
              with
              \begin{equation*}
                  A = \frac{a_1^2}{4} + a_2, \quad B = \frac{a_1a_3}{2} + a_4, \quad C = \frac{a_3^2}{4} + a_6.
              \end{equation*}
              Let $\alpha, \beta, \gamma$ be the distinct roots of the cubic on the right-hand side.
              Applying a linear transformation of $x$, which sends $\alpha, \beta$ to $0, 1$, transforms the equation of $X$ into
              \begin{equation*}
                  y^2 = x(x - 1)(x - \lambda), \quad \lambda = \frac{\gamma - \alpha}{\beta - \alpha}.
              \end{equation*}
              We can plug in $\lambda$ into the $j$-invariant formula.
              Simplifying, we get a rational function of the elementary symmetric polynomials of $\alpha, \beta, \gamma$ (which are just $A, B, C$) with coefficients in $\Q$
              \begin{equation*}
                  j = 2^8 \frac{(A^2 - 3B)^3}{A^2B^2 - 4B^3 -4A^3 + 18ABC - 27C^2}.
              \end{equation*}
              Since $A, B, C$ are polynomials of the $a_i$ with integer coefficients, the $j$-invariant is a rational function of the $a_i$ with coefficients in $\Q$.
          \end{proof}

    \item[\textbf{5.}] Let $X, P_0$ be an elliptic curve having an endormorphism $f : X \to X$ of degree 2.
          \begin{itemize}
              \item[(a)] If we represent $X$ as a 2-1 covering of $\P^1$ by a morphism $\pi : X \to \P^1$ ramified at $P_0$, then as in (4.4), show that there is another morphism $\pi' : X \to \P^1$ and a morphism $g : \P^1 \to \P^1$, also of degree 2, such that $\pi \circ f = g \circ \pi'$.
              \item[(b)] For a suitable choice of coordinates in the two copies of $\P^1$, show that $g$ can be taken to be the morphism $x \mapsto x^2$.
              \item[(c)] Now show that $g$ is branched over two of the branch points of $\pi$, and that $g^{-1}$ of the other two branch points of $\pi$ consists of the four branch points of $\pi'$.
                    Deduce a relation involving the invariant $\lambda$ of $X$.
              \item[(d)] Solving the above, show that there are just three values of $j$ corresponding to elliptic curves with an endomorphism of degree 2, and find the corresponding values of $\lambda$ and $j$.
          \end{itemize}

          \begin{proof} $ $ \vspace{0pt}
              \begin{itemize} [leftmargin=0cm]
                  \item[(a)] The morphism $\pi \circ f$ corresponds to the linear system $|f^*(2P_0)|$.
                        Notice that we can write $f^*(2P_0) = 2f^*(P_0)$, which gives a hint as to how to factor $\pi \circ f$.
                        In particular, we can factor $\pi \circ f$ through the morphism $\pi' : X \to \P^1$ given by the linear system $|f^*(P_0)|$ on $X$, followed by the morphism $g : \P^1 \to \P^1$ given by the linear system $|\pi_*'(f^*P_0)|$ on $\P^1$.

                  \item[(b)] Any degree two morphism $g : \P^1 \to \P^1$ is branched over two points, and we can choose coordinates in the two copies of $\P^1$ so that these two points are $0, \infty$ in the target and source.
                        Then $g$ is given by the morphism $x \mapsto x^2$.
                        Indeed, on some local affine patch, $g$ is given by a degree two polynomial in $x$, say $g(x) = (x - \alpha)(x - \beta)$.
                        Then $g$ is ramified at $x = \frac{1}{2}(\alpha + \beta)$ and $\infty$, and we can choose coordinates so that these two points are $0, \infty$.
                        Similarly, on the target, we can choose coordinates so that the two branch points of $g$ are $0, \infty$.
                        Then $g$ is given by a degree two polynomial with roots at $0, \infty$, which is just $x \mapsto x^2$.

                  \item[(c)] The morphism $\pi \circ f$ is generically four-to-one, and it is branched over the branch points of $\pi$ because any morphism of elliptic curves is unramified by the Riemann-Hurwitz formula.
                        This just says $(\pi \circ f)^{-1}$ of a branch point of $\pi$ consists of two points.
                        Since $g$ is a morphism of degree two ramified at exactly two points, and $\pi'$ is ramified at exactly four points, $g$ must be branched over two of the branch points of $\pi$, and $g^{-1}$ of the other two branch points of $\pi$ consists of the four branch points of $\pi'$.
                        Note that one of the points $g$ is branched over is given by $$g(\pi'(f^{-1}P_0)) = \pi(f(f^{-1}P_0)) = \pi(P_0).$$
                        \filbreak
                        Now, choose coordinates in the two copies of $\P^1$ so that the branch points of $g$ are $0, \infty$, the other two branch points of $\pi$ are $1, \lambda$, and $g$ is given by $x \mapsto x^2$.
                        Then the four branch points of $\pi'$ are $\pm 1, \pm \sqrt{\lambda}$.
                        Consider the linear transformation of $\P^1$ sending $-1, 1, \sqrt{\lambda}$ to $0, 1, \infty$, respectively, which is given by
                        \begin{equation*}
                            \phi(x) = \frac{(1 - \sqrt{\lambda})(x + 1)}{2(x - \sqrt{\lambda})}
                        \end{equation*}
                        Then the branch points of $\pi'$ are sent to $0, 1, \infty, \phi(-\sqrt{\lambda})$, where
                        \begin{equation*}
                            \lambda' := \phi(-\sqrt{\lambda}) = -\frac{(1 - \sqrt{\lambda})^2}{4\sqrt{\lambda}}.
                        \end{equation*}
                        Since $\lambda'$ and $\lambda$ give the same $j$-invariant, they must be related by a fractional linear transformation in $S_3$.
                        In particular, $\lambda'$ is in the orbit of $\lambda$ under the action of $S_3$ on $k - \{0, 1\}$ (4.5).

                  \item[(d)] Solving for $\lambda$ in the above, we get the following three unique values of $\lambda$ modulo the action of $S_3$:
                        \begin{equation*}
                            \lambda = 2, \quad 3 + 2\sqrt{2}, \quad \frac{1 + 3i\sqrt{7}}{2}.
                        \end{equation*}
                        The corresponding values of $j$ are
                        \begin{equation*}
                            j = 2^6 \cdot 3^3, \quad 2^6 \cdot 5^3, \quad -3^3 \cdot 5^3.
                        \end{equation*}
              \end{itemize}
          \end{proof}

    % \item[\textbf{6.}] \begin{itemize}
    %           \item[(a)] Let $X$ be a curve of genus $g$ embedded birationally in $\P^2$ as a curve of degree $d$ with $r$ nodes.
    %                 Generalize the method of (Ex. 2.3) to show that $X$ has $6(g - 1) + 3d$ inflection points.
    %                 A node does not count as an inflection point.
    %                 Assume $\car{k} = 0$.
    %           \item[(b)] Now let $X$ be a curve of genus $g$ embedded as a curve of degree $d$ in $\P^n$, $n \geq 3$, not contained in any $\P^{n - 1}$.
    %                 For each point $P \in X$, there is a hyperplane $H$ containing $P$, such that $P$ counts at least $n$ times in the intersection $H \cap X$.
    %                 This is called an \textit{osculating hyperplane} at $P$.
    %                 It generalizes the notion of tangent line for curves in $\P^2$.
    %                 If $P$ counts at least $n + 1$ times in $H \cap X$, we say $H$ is a \textit{hyperosculating hyperplane}, and that $P$ is a \textit{hyperosculation point}.
    %                 Use Hurwitz's theorem as above, and induction on $n$, to show that $X$ has $n(n + 1)(g - 1) + (n + 1)d$ hyperosculation points.
    %           \item[(c)] If $X$ is an elliptic curve, for any $d \geq 3$, embed $X$ as a curve of degree $d$ in $\P^{d - 1}$, and conclude that $X$ has exactly $d^2$ points of order $d$ in its group law.
    %       \end{itemize}

    %       \begin{proof} $ $ \vspace{0pt}
    %           \begin{itemize} [leftmargin=0cm]
    %               \item[(a)] The rational map $\phi : X \dashrightarrow L$ of (Ex. 2.3a) is defined on the complement of the nodes of $X$.
    %                     By (I, 6.8), $\phi$ extends to a morphism $\phi : X \to L$; to be more specific, if $P$ lies over a node $P'$ of $X$, then $P'$ has two distinct tangent lines, each corresponding to the two points lying over $P'$, and $\phi$ maps $P$ to the point of intersection of $L$ and the tangent line of $X$ at $P'$ corresponding to $P$.
    %                     It follows that the inflection points and the points of intersection of $X$ and $L$ are ramification points of $\phi$ with ramification index 2.

    %               \item[(b)]
    %               \item[(c)] For any point $P \in X$, the linear system $|dP_0|$ is very ample, and it gives an embedding of $X$ as a curve of degree $d$ in $\P^{d - 1}$.
    %                     Plugging in $n = d - 1$ and $g = 1$ into the formula in (b), we get that $X$ has exactly $d(d-2)$ hyperosculation points.
    %           \end{itemize}
    %       \end{proof}

    \item[\textbf{7.}] \textit{The Dual of a Morphism.} Let $X$ and $X'$ be elliptic curves over $k$, with base points $P_0, P_0'$.
          \begin{itemize}
              \item[(a)] If $f : X \to X'$ is any morphism, use (4.11) to show that $f^* : \pic{X'} \to \pic{X}$ induces a homomorphism $\hat{f} : (X', P_0') \to (X, P_0)$.
                    We call this the \textit{dual} of $f$.
              \item[(b)] If $f : X \to X'$ and $g : X' \to X''$ are two morphisms, then $\widehat{g \circ f} = \hat{f} \circ \hat{g}$.
              \item[(c)] Assume $f(P_0) = P_0'$, and let $n = \deg{f}$.
                    Show that if $Q \in X$ is any point, and $f(Q) = Q'$, then $\hat{f}(Q') = n_X(Q)$.
                    Conclude that $f \circ \hat{f} = n_{X'}$ and $\hat{f} \circ f = n_X$.
              \item[(d)] If $f, g : X \to X'$ are two morphisms preserving the base points $P_0, P_0'$, then $\widehat{f + g} = \hat{f} + \hat{g}$.
              \item[(e)] Using (d), show that for any $n \in \Z$, $\hat{n}_X = n_{X}$.
                    Conclude that $\deg n_X = n^2$.
              \item[(f)] Show for any $f$ that $\deg{\hat{f}} = \deg{f}$.
          \end{itemize}

          \begin{proof} Let the pair $(X, \fP)$ (resp. $(X', \fP')$) be the Jacobian of $(X, P_0)$ (resp. $(X', P_0')$).
              If $\fC$ is any category, and $T$ is any object in $\fC$, then denote by $h_T : \fC^\text{op} \to \cSet$ the contravariant Hom functor defined by $U \mapsto \Hom_\fC(U, T)$.
              Recall that $X$ and $\fP \in \pic^\circ(X/X)$ satisfies the following universal property: if $T$ is any scheme of finite type over $k$, and for any $\fM \in \pic^\circ(X/T)$, there is a unique morphism $f : T \to X$ such that $f^*\fP \cong \fM$ in $\pic^\circ(X/T)$.
              In other words, $(X, \fP)$ represents the functor $\pic^\circ(X/-)$.
              Lastly, since $(X, \fP)$ is a group object in the category of schemes of finite type over $k$, the set $h_X(T)$ has a group structure induced by $\pic^\circ(X/T)$.

              \begin{itemize} [leftmargin=0cm]
                  \item[(a)] Any morphism $f : X \to X'$ defines a natural transformation $f^* : \pic^\circ(X' / -) \to \pic^\circ(X / -)$ of group functors in the following way (4.10.2).

                        If $T$ is any scheme of finite type over $k$, then $f' = f \times 1_T : X' \times T \to X \times T$ induces a group homomorphism $f^* : \pic(X' \times T) \to \pic(X \times T)$.
                        Let $p : X \times T \to T$ and $p' : X' \times T \to T$ be the projection maps.
                        The image of $f^*$ lies in $\pic^\circ(X \times T)$, and $p = p' \circ f'$ implies $f^* p'^* \pic(T) = p^* \pic(T)$.
                        Hence, we have a well-defined group homomomorphism $f^*_T : \pic^\circ(X'/T) \to \pic^\circ(X/T)$ for each scheme $T$ of finite type over $k$.

                        It remains to show $f^*$ is functorial.
                        Let $g : S \to T$ be any morphism of schemes of finite type over $k$.
                        We can apply $\pic(-)$ to the following commutative diagram
                        \[ \begin{tikzcd}
                                X \times S \arrow[d, "1_X \times g"'] \arrow[r, "f \times 1_S"] & X' \times S \arrow[d, "1_{X'} \times g"] \\
                                X \times T \arrow[r, "f \times 1_T"]                 & X' \times T
                            \end{tikzcd}, \]
                        which gives a commutative diagram of Picard groups with arrows given by taking inverse images of invertible sheaves.

                        Now, we identify $h_X$ and $h_{X'}$ with $\pic^\circ(X/-)$ and $\pic^\circ(X'/-)$, respectively.
                        Yoneda's lemma says there exists a natural bijection between the set of morphisms $X' \to X$ and the set of natural transformations of functors $h_{X'} \to h_X$.
                        Hence, there exists a unique morphism $\hat{f} \in \Hom(X', X)$ such that $f^* \in \Hom(h_{X'}, h_X)$ is given by composition with $\hat{f}$.
                        The invertible sheaves $\fO_X$ and $\fO_{X'}$ correspond to $P_0 \in X$ and $P_0' \in X'$, respectively, and structure sheaves pullback, so we conclude that $\hat{f}(P_0') = P_0$.

                  \item[(b)] This is a consequence of the uniqueness property of $\hat{f}$ in $\Hom(X', X)$, and the analogous fact for inverse image of sheaves, i.e., $(g \circ f)^* = f^* \circ g^*$.

                  \item[(c)] The points $Q' \in X'$ and $\hat{f}(Q') \in X$ correspond to $\fO_{X'}(Q' - P_0')$ in $\pic^\circ(X'/k)$ and $f^*\fO_{X'}(Q' - P_0')$ in $\pic^\circ(X/k)$, respectively.
                        Thus, it is enough to show the linear equivalence of the following divisors
                        \begin{equation*}
                            f^*(Q' - P_0') \sim n_X(Q) - P_0.
                        \end{equation*}
                        By the hypothesis, $Q' - P_0' = f_*(Q - P_0)$, which by (Ex. 2.6) implies
                        \begin{equation*}
                            f^*(Q' - P_0') = n(Q - P_0).
                        \end{equation*}
                        The right-hand side is linearly equivalent to $n_X(Q) - P_0$ because the group law on $X$ is induced by that of $\pic^\circ(X/k)$, so we conclude that $\hat{f} \circ f = n_X$.
                        For the second corollary, observe that $f$ and multiplication by $n$ commutes since $f$ is a morphism of proper group schemes.
                        Hence, we have
                        \begin{equation*}
                            (f \circ \hat{f})(Q') = f(n_X(Q)) = n_{X'}(f(Q)) = n_{X'}(Q').
                        \end{equation*}

                  \item[(d)] It is enough to show for any $\fL' \in \pic^\circ(X')$, that
                        \begin{equation*}
                            (f + g)^*\fL' \cong f^*\fL' \otimes g^*\fL'.
                        \end{equation*}
                        Let $Q'$ be the closed point in $X'$ such that $\fL' \cong \fO_{X'}(Q' - P_0')$.
                        There exists a bijection
                        \begin{align*}
                            \Hom_{\ckSch{k}}(X, X') & \xrightarrow{\sim} \pic^\circ(X'/X)       \\
                            f : X \to X'            & \mapsto \fL_f := (f \times 1_{X'})^*\fP',
                        \end{align*}
                        where the group law on the left-hand side is induced by $\pic^\circ(X'/X)$, i.e., $\fL_{f + g} = \fL_f \otimes \fL_g$ is a representative of $f + g$ in $\pic(X \times X')$ for any $f, g : X \to X'$.
                        Let $\sigma : X \to X \times X'$ be the section $x \mapsto (x, Q')$.
                        If we assume $\fP' = \fO_{X' \times X'}(\Delta - X' \times P_0')$ as in (4.11), then $\sigma^* \fL_f \cong f^*\fO_{X'}(Q')$, so replacing $\fP'$ by $\fP' \otimes p_2^*\fO_{X'}(-P_0)$, we get $\sigma^*\fL_f \cong f^*\fL'$.
                        Hence,
                        \begin{align*}
                            (f + g)^* \fL' \cong \sigma^* \fL_{f + g} \cong \sigma^*\fL_f \otimes \sigma^* \fL_g \cong f^*\fL' \otimes g^*\fL'.
                        \end{align*}

                  \item[(e)] Proceeding by induction on $n$ reduces the statement to the case $n = 2$.
                        The dual of the identity is the identity.
                        Hence, $\hat{2}_X = \widehat{1_X + 1_X} = 1_X + 1_X = 2_X$ by (d).
                        According to (c), $\hat{n}_X \circ n_X$ is multiplication by $\deg{n_X}$.
                        Since $\hat{n}_X \circ n_X = n^2_X$, we conclude that $\deg{n_X} = n^2$.

                  \item[(f)] By (c) and (e), $\deg \hat{f} \circ f = n^2$.
                        The degree of morphisms of curves is multiplicative with respect to composition.
                        Hence, $\deg{\hat{f}} = n$.
              \end{itemize}
          \end{proof}

    \item[\textbf{9.}] We say two elliptic curves $X, X'$ are \textit{isogenous} if there is a finite morphism $f : X \to X'$.

          \begin{itemize}
              \item[(a)] Show that isogeny is an equivalence relation.
              \item[(b)] For any elliptic curve $X$, show that the set of elliptic curves $X'$ isogenous to $X$, up to isomorphism is countable.
          \end{itemize}

          \begin{proof} $ $ \vspace{0pt}
              \begin{itemize} [leftmargin=0cm]
                  \item[(a)] The identity morphism is trivially finite, so any elliptic curve is isogenous to itself.
                        A composition of finite morphisms is finite, so isogeny satisfies transitivty.
                        Lastly, (Ex. 4.7) implies isogeny is reflexive.
                        Hence, isogeny is an equivalence relation.

                  \item[(b)] Since $\ker{f}$ is a closed subgroup scheme of $X$, it is a finite set of points, all of which are torsion under the group law of $X$.
                        Every torsion point of $X$ lies in the kernel of a multiplication map $n_X : X \to X$ for some $n > 0$.
                        By (Ex. 4.7), $n_X$ is an étale covering of degree $n^2$, which implies $X$ has exactly $n^2$ points of order $n$.
                        Hence, there are countably many torsion points on $X$.
                        The set of finite subsets of a countable set is countable, so there are only countably many candidates for $\ker{f}$.
                        Since $X'$ is uniquely determined by $X$ and $\ker{f}$, the set of elliptic curves isogenous to $X$ is countable.
              \end{itemize}
          \end{proof}

    \item[\textbf{11.}] Let $X$ be an elliptic curve over $\C$, defined by the elliptic functions with periods $1, \tau$.
          Let $R$ be the the ring of endormorphisms of $X$.

          \begin{itemize}
              \item[(a)] If $f \in R$ is a nonzero endormorphism correspondending to complex multiplication by $\alpha$ as in (4.18), show that $\deg{f} = |\alpha|^2$.
              \item[(b)] If $f \in R$ corresponds to $\alpha \in \C$ again, show that the dual $\hat{f}$ of (Ex. 4.7) corresponds to the complex conjugate $\bar{\alpha}$ of $\alpha$.
              \item[(c)] If $\tau \in \Q(\sqrt{-d})$ happens to be integral over $\Z$, show that $R = \Z[\tau]$.
          \end{itemize}

          \begin{proof} $ $ \vspace{0pt}
              \begin{itemize} [leftmargin=0cm]
                  \item[(a)] It is not hard to see $\deg{f}$ is the size of the quotient $|\Lambda / \alpha \Lambda|$.
                        Equivalently, it is the ratio between the areas of the  parallelograms with vertices $\alpha, \alpha \tau$, and $1, \tau$.
                        But $\alpha$ just scales the area by $|\alpha|^2$.
                        Hence, $\deg{f} = |\alpha|^2$.

                  \item[(b)] By (Ex. 7c), $\hat{f} \circ f$ is multiplication by $|\alpha|^2$, but $f$ is multiplication by $\alpha$, and $\bar{\alpha}$ is the unique element in $\C$ such that $\alpha \bar{\alpha} = |\alpha|^2$.
                        Hence, $\hat{f}$ is multiplication by $\bar{\alpha}$.

                  \item[(c)] The ring of integers of $\Q(\sqrt{-d})$ is given by
                        \begin{equation*}
                            \begin{cases*}
                                \Z[\sqrt{-d}]                          & $d \equiv 1, 2 \mod{4}$ \\
                                \Z\bigg[\cfrac{1 + \sqrt{-d}}{2}\bigg] & $d \equiv 3 \mod{4}$
                            \end{cases*}.
                        \end{equation*}
                        By (4.19), if $\tau = r + s \sqrt{-d}$ for some $r, s \in \Q$, the endomorphism ring $R$ is
                        \begin{equation*}
                            R = \{ a + b \tau \mid a, b \in \Z, \quad 2br, b(r^2 + ds^2) \in \Z \}.
                        \end{equation*}
                        Suppose $\tau$ is integral over $\Z$.
                        If $d \equiv 1, 2 \mod{4}$, then $r, s \in \Z$, so $a + b \tau \in R$ for all $a, b \in \Z$.
                        If $d \equiv 3 \mod{4}$, then we can write $\tau = (r' + s'\sqrt{-d}) / 2$ with $r', s' \in \Z$ such that $r' \equiv s' \mod{2}$.
                        For any $b \in \Z$, the first condition $2br = br' \in \Z$ is easily satisfied.
                        Also, we have
                        $$r^2 + ds^2 = \frac{r'^2 + ds'^2}{4} \in \Z$$ for all $b \in \Z$ since $r'^2, s'^2 \equiv 0, 1 \mod{4}$.
                        Hence, $R = \Z[\tau]$.
              \end{itemize}
          \end{proof}

    \item[\textbf{12.}] Again let $X$ be an elliptic curve over $\C$ determined by the elliptic functions with periods $1, \tau$, and assume that $\tau$ lies in the region $G$ of (4.15B).

          \begin{itemize}
              \item[(a)] If $X$ has any autmorphisms leaving $P_0$ fixed other than $\pm 1$, show that either $\tau = i$ or $\tau = \omega$, as in (4.20.1) and (4.20.2).
                    This gives another proof of the fact (4.7) that there are only two curves, up to isomorphism, having automorphisms other than $\pm 1$.
              \item[(b)] Now show that there are exactly three values of $\tau$ for which $X$ admits an endomorphism of degree 2.
                    Can you match these with the three values of $j$ determined in (Ex. 4.5)?
          \end{itemize}

          \begin{proof} $ $ \vspace{0pt}
              \begin{itemize} [leftmargin=0cm]
                  \item[(a)] By (Ex. 4.11), $\tau$ must be a root of unity in $G$.
                        The only roots of unity in $G$ are $i$ and $\omega$.

                  \item[(b)] By (Ex. 4.11), we need to find the values of $\tau$ such that $\Lambda$ contains an element of norm 2.
                        In particular, if we write $m + n\tau \in \Lambda$, then we need to determine when
                        \begin{equation*}
                            m^2 + n^2|\tau|^2 + 2mn\Re(\tau) = 4
                        \end{equation*}
                        admits a solution for some $m, n \in \Z$.
                        Using the bounds on $|\tau|$ and $\Re(\tau)$ as specified in (4.15B), the only possible values of $\tau$ are
                        \begin{equation*}
                            \tau = i, \quad i\sqrt{2}, \quad -\frac{1}{2} + i\frac{\sqrt{7}}{2}.
                        \end{equation*}
                        The corresponding values of $j$ are
                        \begin{equation*}
                            j = 2^6 \cdot 3^3, \quad 2^6 \cdot 5^3, \quad - 3^3 \cdot 5^3.
                        \end{equation*}
              \end{itemize}
          \end{proof}

    \item[\textbf{13.}] If $p = 13$, there is just one value of $j$ for which the Hasse invariant of the corresponding curve is 0.
          Find it.

          \begin{proof}
              The polynomial $h_p(\lambda)$ of (4.22) for $p = 13$ is
              \begin{align*}
                  h_{13}(\lambda) = \lambda^6 + 10\lambda^5 + 4\lambda^4 + 10\lambda^3 + 4\lambda^2 + 10\lambda + 1,
              \end{align*}
              and an elliptic curve has Hasse invariant 0 if and only if $h_{13}(\lambda) = 0$.
              On the otherhand, consider the $j$-invariant
              \begin{equation*}
                  j = 2^8 \frac{(\lambda^2 - \lambda + 1)^3}{\lambda^2(\lambda - 1)^2}.
              \end{equation*}
              Expanding above modulo 13 gives
              \begin{equation*}
                  \lambda^6 + 10\lambda^5 + (10j + 6)\lambda^4 + (6j + 6)\lambda^3 + (10j + 6)\lambda^2 + 10\lambda + 1 = 0.
              \end{equation*}
              Matching coefficients with $h_{13}(\lambda)$ and solving for $j$ gives $j \equiv 5 \mod{13}$.
          \end{proof}

    \item[\textbf{14.}] The Fermat curve $X : x^3 + y^3 = z^3$ gives a nonsingular curve in characteristic $p$ for every $p \neq 3$.
          Determine the set $\gP = \{p \neq 3 \mid \text{$X_{(p)}$ has Hasse invariant 0} \}$, and observe (modulo Dirichlet's theorem) that it is a set of primes of density $1 / 2$.

          \begin{proof}
              Using (4.21) with $f(x, y, z) = x^3 + y^3 - z^3$, we find that the coefficient of $(xyz)^{p - 1}$ in $f^{p - 1}$ is zero if and only if $p \equiv 2 \mod{3}$.
              By Dirichlet's theorem, the density of $\goth{P}$ is $1 / \phi(3) = 1 / 2$.
          \end{proof}

    \item[\textbf{15.}] Let $X$ be an elliptic curve over a field $k$ of characteristic $p$.
          Let $F' : X_p \to X$ be the $k$-linear Frobenius morphism (2.4.1).
          Use (4.10.7) to show that the dual morphism $\hat{F}' : X \to X_p$ is separable if and only if the Hasse invariant of $X$ is 1.
          Now use (Ex. 4.7) to show that if the Hasse invariant is 1, then the subgroup of points of order $p$ on $X$ is isomorphic to $\Z / p$; if the Hasse invariant is 0, then it is 0.

          \begin{proof}
              By (III, Ex. 8.1) and (4.10.7), if $f : X \to X'$ is any morphism of elliptic curves over $k$, then the induced map of cohomology $f^* : \cH^1(X, \fO_{X'}) \to \cH^1(X, \fO_X)$ is precisely the map of tangent spaces at 0 of the dual map $\hat{f} : X' \to X$.
              By the Riemann-Hurwitz formula, any separable morphism of elliptic curves is not ramified.
              Thus, a morphism of elliptic curves $f : X \to X'$ is inseparable if and only if $f$ is ramified at a point, in which case $f$ is ramified everywhere by the group structure of elliptic curves.
              Also, any morphism of curves is ramified at a point if and only if the map of tangent spaces is not an isomorphism.
              Hence, $\hat{F}'$ is separable if and only if the map of tangent spaces $F'^* : \cH^1(X, \fO_X) \to \cH^1(X_p, \fO_{X_p})$ is not the zero map if and only if the Hasse invariant of $X$ is 1.

              The subgroup of points of order $p$ on $X$ is the kernel of the multiplication by $p$ morphism $p_X : X \to X$.
              Since the $k$-linear Frobenius morphism $F'$ has degree $p$, we have $p_X = F' \circ \hat{F}'$ by (Ex. 4.7c).
              Also, $F'$ is a homeomorphism on the underlying topological spaces, i.e., $F'$ is a one-to-one map.
              If the Hasse invariant of $X$ is 1, then $\hat{F}'$ is a separable morphism of degree $p$, so $\hat{F}'$ is a $p$-to-one map.
              Hence, $p_X$ is a $p$-to-one map, and $(\ker{p_X})[k] \cong \Z / p$.
              If the Hasse invariant of $X$ is 0, then $p_X = F' \circ \hat{F}'$ is a purely inseparable morphism of curves.
              Hence, $p_X$ is a one-to-one map, and $(\ker{p_X})[k] = 0$.
          \end{proof}

    \item[\textbf{16.}] Again let $X$ be an elliptic curve over $k$ of characteristic $p$, and suppose $X$ is defined over the field $\F_q$ of $q = p^r$ elements, i.e., $X \subseteq \P^2$ can be defined by an equation with coefficients in $\F_q$.
          Assume also that $X$ has a rational point over $\F_q$.
          Let $F' : X_q \to X$ be the $k$-linear Frobenius with respect to $q$.

          \begin{itemize}
              \item[(a)] Show that $X_q \cong X$ as schemes over $k$, and that under this identification, $F' : X_q \to X$ is the map obtained by the $q$th-power map on the coordinates of points of $X$, embedded in $\P^2$.
              \item[(b)] Show that $1_X - F'$ is a separable morphism and its kernel is just the set $X(\F_q)$ of points of $X$ with coordinates in $\F_{q}$.
              \item[(c)] Using (Ex. 4.7), show that $F' + \hat{F}' = a_X$ for some integer $a$, and that $N = q - a + 1$, where $N = \# X(\F_q)$.
              \item[(d)] Use the fact that $\deg(m + nF') > 0$ for all $m, n \in \Z$ to show that $|a| \leq 2\sqrt{q}$.
                    This is Hasse's proof of the analogue of the Riemann hypothesis for elliptic curves.
              \item[(e)] Now assume $q = p$, and show that the Hasse invariant of $X$ is 0 if and only if $a \equiv 0 \mod{p}$.
                    Conclude that for $p \geq 5$ that $X$ has Hasse invariant 0 if and only if $N = p + 1$.
          \end{itemize}

          \begin{proof} $ $ \vspace{0pt}
              \begin{itemize} [leftmargin=0cm]
                  \item[(a)] Suppose $X$ is locally isomorphic to $\spec{A}$, where $A = k[x, y]/ (f)$ and $f(x, y)$ is some degree three polynomial with coefficients in $\F_q \subset k$.
                        Then $X_p$ is locally isomorphic to $\spec{(k \otimes_q A)}$, where the copy of $k$ here has the $q$th-power Frobenius for the structure morphism as a $k$-algebra.
                        The image of $f$ under the identification
                        \begin{align*}
                            k \otimes_q k[x, y] \cong k[x, y] \\
                            (b \cdot 1) x \mapsto b^qx,
                        \end{align*}
                        gives a defining equation $f^{(q)}$ for $X_q$.
                        Notice that $f^{(q)}$ is obtained from $f$ by raising each coefficient of $f$ to the $q$th-power.
                        Since each coefficient of $f$ is in $\F_q$, we see that $f^{(q)}$ and $f$ are actually equal in $k[x, y]$.
                        Hence, $X_q \cong X$ as schemes over $k$.

                        Under this indentification, $F' : X_q \to X$ is obtained from the restriction of the Frobenius on $\P^2$.
                        The $k$-linear Frobenius morphism on an affine coordinate ring of $\P^2$ can be written as
                        \begin{align*}
                            F^* : k[x, y] & \to k[x, y]          \\
                            g(x, y)       & \mapsto g(x^q, y^q),
                        \end{align*}
                        which is exactly the ring homomorphism that gives the $q$th-power map on the coordinates.

                  \item[(b)] If $t \in \fO_{P, X}$ is a local parameter of $X$, then $(1_X - F')^*t = t - t^p$.
                        Since $\dd(t - t^p) = \dd t$, the map $1_X - F'$ is ramified nowhere.
                        Hence, $1_X - F'$ is a separable morphism.
                        Under the embedding $X \subset \P^2$, the kernel of $1_X - F'$ consists of the points fixed by the map on $X$ obtained by the $q$th-power map on the coordinates of points of $X$.
                        Since $\F_q$ is precisely the fixed field of $k$ under the $q$th-power Frobenius, the kernel of $1_X - F'$ consists of the points of $X \subset \P^2$ with coordinates in $\F_q$.

                  \item[(c)] Recall from (Ex. 4.7), (Ex. 4.15), and (Ex. 4.16b) that
                        \begin{equation*}
                            (1_X-F')\circ\widehat{(1_X-F')}=N_X, \quad F'\circ\hat F'=q_X.
                        \end{equation*}
                        Expanding the left-hand side of the first identity gives
                        \begin{equation*}
                            (1_X-F')\circ\widehat{(1_X-F')} = 1_X-(F'+\hat F') + q_X.
                        \end{equation*}
                        Hence, we have
                        \begin{equation*}
                            F' + \hat F' = q_X + 1_X - N_X.
                        \end{equation*}

                  \item[(d)] For all $m, n \in \Z$, we have the inequality
                        \begin{equation*}
                            m^2 + n^2 q - mna = \deg(m - nF') > 0.
                        \end{equation*}
                        Rearranging gives
                        \begin{equation*}
                            |a| < \bigg| \frac{m}{n} + \frac{n}{m}q\bigg|
                        \end{equation*}
                        The real-valued function $f(t) = t + q / t$ has a local extremas at $(\sqrt{q}, \pm 2 \sqrt{q})$.
                        The rational numbers are dense in the real numbers, so we conclude that $|a| \leq 2 \sqrt{q}$.
                        I wonder if it is every the case $a = 2\sqrt{q}$ when $\sqrt{q}$ is a perfect square.

                  \item[(e)] By (Ex. 4.15), the Hasse invariant of $X$ is 0 if and only if $\hat{F}'$ is inseparable.
                        If $t \in \fO_{P, X}$ is a local parameter, then
                        \begin{equation*}
                            (F' + \hat{F}')^*t = t^p + \hat{F}'^* t,
                        \end{equation*}
                        so $\hat{F}'$ is inseparable if and only if $F' + \hat{F}'$ is inseparable.
                        The conclusion follows from the fact that $a_X$ is inseparable if and only if $a \equiv 0 \mod{p}$.
                        For $p \geq 5$, $a \equiv 0 \mod{p}$ and $|a| \leq 2 \sqrt{p}$ can be both satisfied if and only if $a = 0$.
                        What about the cases for $p < 5$?
              \end{itemize}
          \end{proof}

    \item[\textbf{17.}] Let $X$ be the curve $y^2 + y = x^3 - x$ of (4.23.8).

          \begin{itemize}
              \item[(a)] If $Q = (a, b)$ is a point on the curve, compute the coordinates of the point $P + Q$, where $P = (0, 0)$, as a function of $a, b$.
                    Use this formula to find the coordinates $nP$, $n = 1, 2, \dots, 10$.
              \item[(b)] This equation defines a nonsingular curve over $\F_q$ for all $p \neq 37$.
          \end{itemize}

          \begin{proof} $ $ \vspace{0pt}
              \begin{itemize} [leftmargin=0cm]
                  \item[(a)] $P + Q = \bigg(-\cfrac{a + b}{a^2}, \cfrac{b(a + b)}{a^3} - 1\bigg), \quad -P = (0, -1), \quad 2P = (1, 0)$

                        \begin{tikzpicture}
                            \begin{axis}[
                                    width=16cm,
                                    height=12cm,
                                    axis lines=middle,
                                    xlabel={$x$},
                                    ylabel={$y$},
                                    xmin=-1.2,
                                    xmax=2.1,
                                    ymin=-4,
                                    ymax=3.5,
                                    samples=300,
                                    clip=false,
                                ]

                                % Upper branch
                                \addplot[
                                    blue,
                                    thick,
                                    domain=-1.107159871689:0.269594436405
                                ]
                                {-1/2 + sqrt(x^3 - x + 1/4)};

                                \addplot[
                                    blue,
                                    thick,
                                    domain=0.837565435283:2.1,
                                ]
                                {-1/2 + sqrt(x^3 - x + 1/4)};

                                % Lower branch
                                \addplot[
                                    blue,
                                    thick,
                                    domain=-1.107159871689:0.269594436405
                                ]
                                {-1/2 - sqrt(x^3 - x + 1/4)};

                                \addplot[
                                    blue,
                                    thick,
                                    domain=0.837565435283:2.1,
                                ]
                                {-1/2 - sqrt(x^3 - x + 1/4)};


                                \addplot[
                                    only marks,
                                    mark=*,
                                    mark size=2pt,
                                    red,
                                ]
                                coordinates {
                                        (0,-1)
                                        (0, 0)
                                        (1, 0)
                                        (-1,-1)
                                        (2,-3)
                                        (0.25,-0.625)
                                        % (6,14)
                                        (-5/9,8/27)
                                        (21/25,-69/125)
                                        (-20/49,-435/343)
                                        % (161/16,-2065/64)
                                    };

                                % Labels
                                \node[anchor=north west] at (axis cs:0,-1) {$-P$};
                                \node[anchor=south west] at (axis cs:0,0) {$P$};
                                \node[anchor=south] at (axis cs:1,0) {$2P$};
                                \node[anchor=north] at (axis cs:-1,-1) {$3P$};
                                \node[anchor=north] at (axis cs:2,-3) {$4P$};
                                \node[anchor=west] at (axis cs:0.25,-0.625) {$5P$};
                                % \node[anchor=west] at (axis cs:6,14) {$(6,14)$};
                                \node[anchor=south] at (axis cs:-5/9,8/27) {$7P$};
                                \node[anchor=east] at (axis cs:21/25,-69/125) {$8P$};
                                \node[anchor=north] at (axis cs:-20/49,-435/343) {$9P$};
                                % \node[anchor=west] at (axis cs:161/16,-2065/64) {$\left(\frac{161}{16},-\frac{2065}{64}\right)$};
                            \end{axis}
                        \end{tikzpicture}

                  \item[(b)] Let $X$ be a cubic curve in $\P^2$ given by an equation of the form
                        \begin{equation*}
                            y^2 + a_1xy + a_3y = x^3 + a_2x^2 + a_4x + a_6, \quad a_i \in k.
                        \end{equation*}
                        Define the following quantities
                        \begin{align*}
                            b_2 & = a_1^2 + 4a_2                                            \\
                            b_4 & = 2a_4 + a_1 a_3                                          \\
                            b_6 & = a_3^2 + 4a_6                                            \\
                            b_8 & = a_1^2 a_6 + 4a_2 a_6 - a_1 a_3 a_4 + a_2 a_3^2 - a_4^2.
                        \end{align*}
                        The discriminant of $X$ is given by
                        \begin{equation*}
                            \Delta = -b_2^2 b_8 - 8 b_4^3 - 27 b_6^2 + 9 b_2 b_4 b_6,
                        \end{equation*}
                        and $X$ defines a nonsingular curve over $\F_q$ if and only if $\Delta \neq 0$ in $\F_q$.
                        For our elliptic curve, the quantities $b_i$ are
                        \begin{equation*}
                            b_2 = 0, \quad b_4 = -2, \quad b_6 = 1, \quad b_8 = -1,
                        \end{equation*}
                        so the determinant of $X$ is $\Delta = 37$.
                        Hence, $X$ defines a nonsingular curve over $\F_q$ for all $p \neq 37$.
              \end{itemize}
          \end{proof}

    \item[\textbf{20.}] Let $X$ be an elliptic curve over a field $k$ of characteristic $p > 0$, and let $R = \End(X, P_0)$ be its ring of endormorphisms.

          \begin{itemize}
              \item[(a)] Let $X_p$ be the curve over $k$ defined by changing the $k$-structure of $X$ (2.4.1).
                    Show that $j(X_p) = j(X)^{1/p}$.
                    Thus, $X \cong X_p$ over $k$ if and only if $j \in \F_p$.

              \item[(b)] Show that $p_X$ in $R$ factors into a product of $\pi \hat{\pi}$ of two elements of degree $p$ if and only if $X \cong X_p$.
                    In this case, the Hasse invariant of $X$ is 0 if and only if $\pi$ and $\hat{\pi}$ are associates in $R$.

              \item[(c)] If $\Hasse(X) = 0$ show in any case $j \in \F_{p^2}$.

              \item[(d)] For any $f \in R$, there is an induced map $f^* : \cH^1(\fO_X) \to \cH^1(\fO_X)$.
                    This must be multiplication by an element $\lambda_f \in k$.
                    So we obtain a ring homomorphism $\varphi : R \to k$ by sending $f$ to $\lambda_f$.
                    Show that any $f \in R$ commutes with the (nonlinear) Frobenius morphism $F : X \to X$, and conclude that if $\Hasse(X) \neq 0$, then the image of $\varphi$ is $\F_p$.
                    Therefore, $R$ contains a prime ideal $\goth{p}$ with $R / \goth{p} \cong \F_p$.
          \end{itemize}

          \begin{proof} $ $ \vspace{0pt}
              \begin{itemize} [leftmargin=0cm]
                  \item[(a)] Some of the statements proved in (Ex. 4.16) hold here as well.
                        In particular, under the idenficiation of the $k$-linear Frobenius morphism $F^* : k[x, y] \to k[x, y]$ defined by $x, y \mapsto x^p , y^q$, let $f(x, y)$ be a cubic polynomal that locally defines $X$ in $\P^2$.
                        The image of $f(x, y)$ under $F^*$ is $f(x^p, x^p) = 0$, so if we write
                        \begin{equation*}
                            f(x, y) = y^2 - x(x - 1)(x - \lambda),
                        \end{equation*}
                        then
                        \begin{equation*}
                            F^*f(x, y) = y^{2p} - x^p(x^p - 1)(x^p - \lambda).
                        \end{equation*}
                        Since $k$ is algebraically closed, there exists a unique $p$th root $\lambda^{1/p}$ of $\lambda$.
                        Thus, we can factorize
                        \begin{equation*}
                            y^{2p} - x^p(x^p - 1)(x^p - \lambda) = (y^2 - x(x - 1)(x - \lambda^{1/p}))^p.
                        \end{equation*}
                        The radical of the right-hand side gives the defining equation of $X_p$.
                        Plugging in $\lambda^{1/p}$ gives $j(X_p) = j(X)^{1/p}$.

                  \item[(b)] If $X$ admits a degree $p$ endomorphism $\pi : X \to X$, then $X \cong X_p$ by (2.5).
                        Conversely, if $X_p \cong X$, then $p_X$ factors as $F' \circ \hat{F}'$, where $F' : X_p \to X$ is the $k$-linear Frobenius morphism (Ex. 4.7).
                        By (Ex. 4.15), the Hasse invariant of $X$ is 0 if and only if $\hat{F}'$ is an inseparable morphism.
                        The Frobenius is the unique purely inseparable morphism of degree $p$ (2.5), so if $X_p \cong X$, then $F'$ and $\hat{F}'$ must differ by an automorphism of $X$, i.e., they are associates in $R$.
                        If $F'$ and $\hat{F}'$ are associates, then $\hat{F}'$ is an inseparable morphism since $F'$ is inseparable.
                        Hence, the Hasse invariant of $X$ is 0.

                  \item[(c)] If $\Hasse(X) = 0$, then $\hat{F}' : X \to X_p$ is a purely inseparable morphism of degree $p$.
                        By (2.5), $X$ and the $p$th Frobenius twist $(X_{p})_p \cong X_{p^2}$ of $X_p$ are isomorphic as abstract $k$-schemes.
                        Hence, $j(X) = j(X)^{1/p^2}$ by (Ex. 4.20a).

                  \item[(d)] Clearly $f$ and $F$ commute on the underlying topological spaces.
                        The statement holds for any homomorphism between $k$-algebras, so it is immediate that $f$ and $F$ commute at the level of sheaves as well.
                        If $\Hasse(X) \neq 0$, then the $p$-linear map of cohomology $F^* : \cH^1(X, \fO_X) \to \cH^1(X, \fO_X)$ is non-zero, and commutes with $f^* : \cH^1(X, \fO_X) \to \cH^1(X, \fO_X)$.
                        In particular, $F^*(1) \neq 0$, and $(f \circ F)^*(1) = (F \circ f)^*(1)$.
                        Expanding both sides gives
                        \begin{align*}
                            (f \circ F)^*(1) & = F^*(f^*(1)) = F^*(\lambda_f \cdot 1) = \lambda_f^p \cdot F^*(1) \\
                            (F \circ f)^*(1) & = f^*(F^*(1)) = \lambda_f \cdot F^*(1).
                        \end{align*}
                        Since $F^*(1)$ is non-zero, we have $\lambda_f^p = \lambda_f$.
                        Hence, $\lambda_f$ is in $\F_p$.
              \end{itemize}
          \end{proof}
\end{enumerate}

\end{document}
